\chapter{Notation}
\label{ch:notation}

Bold upright capitals denote second-order tensors and matrices, bold lowercase
denotes vectors, and blackboard bold denotes fourth-order tensors.  Bold Greek
follows the same convention through the \code{boldtensors} package: $\Cauchy$
is the Cauchy stress tensor, $\sigma$ a scalar.

\section*{Continuum quantities}

\begin{tabular}{@{}ll@{}}
\toprule
Symbol & Meaning \\
\midrule
$\vect{X}$, $\vect{x}$        & material and spatial position \\
$\disp$                       & displacement field, $\vect{x} = \vect{X} + \disp$ \\
$\Fdef = \Ident + \Grad\disp$ & deformation gradient \\
$J = \det\Fdef$               & Jacobian of the deformation \\
$\Cright = \Fdef^{\mathsf{T}}\Fdef$ & right Cauchy--Green tensor \\
$\Piola$                      & first Piola--Kirchhoff stress \\
$\Cauchy$, $\Kirch = J\Cauchy$ & Cauchy and Kirchhoff stress \\
$\Psi$                        & strain-energy density per unit reference volume \\
$\fourth{A} = \partial\Piola/\partial\Grad\disp$ & material tangent \\
$\rho_0$                      & reference mass density \\
\bottomrule
\end{tabular}

\section*{Discrete quantities}

\begin{tabular}{@{}ll@{}}
\toprule
Symbol & Meaning \\
\midrule
$n$                     & number of unknown (free) degrees of freedom \\
$\Uvec, \Vvec, \Avec$   & nodal displacement, velocity, acceleration \\
$\Rvec$                 & residual (out-of-balance force) \\
$\Kmat = \partial\Rvec/\partial\Uvec$ & tangent stiffness \\
$\Mmat$                 & mass matrix \\
$\Keff$                 & effective operator of an implicit step \\
$c_M$                   & mass coefficient of $\Keff$, Section~\ref{sec:newmark} \\
$\Dmat = \diag(\Kmat)$  & diagonal of an operator \\
$\Prec$                 & preconditioner, $\Prec \approx \Kmat^{-1}$ \\
$\Pmat$, $\Rrest$       & prolongation and restriction \\
$\Bnull$                & near-nullspace, Section~\ref{sec:nullspace} \\
$\eta_k$                & forcing term at Newton iteration $k$, Chapter~\ref{ch:inexact} \\
$\lambda_{\max}$        & largest eigenvalue of $\Dmat^{-1}\Kmat$ \\
\bottomrule
\end{tabular}

\section*{Index conventions}

Subscript $k$ indexes nonlinear (Newton) iterations, subscript $i$ or $j$
indexes linear (Krylov) iterations, subscript $n$ indexes time steps, and
superscript $\ell$ indexes multigrid levels.  A quantity written without a time
subscript is understood at $t_{n+1}$, the step being solved.

Norms are Euclidean unless a subscript says otherwise; $\norm{\vect{r}}_{\Prec}
= \sqrt{\vect{r}^{\mathsf{T}}\Prec\vect{r}}$ denotes the norm induced by a
preconditioner, which is the quantity preconditioned conjugate gradients
actually monitors (Section~\ref{sec:pcg}).

\section*{Free and constrained degrees of freedom}

Carina imposes Dirichlet conditions by \emph{elimination}, never by penalty.
The global degree-of-freedom set partitions into unknowns and prescribed
values, and every operator named above is understood as the unknown--unknown
block unless stated otherwise.  Chapter~\ref{ch:formulation} makes this precise
and Section~\ref{sec:elimination} explains why the choice matters for every
Krylov method in this manual.
